\documentclass[11pt]{article}
\usepackage[T1]{fontenc}
\usepackage{lmodern,amsmath,amssymb,amsthm,microtype}
\usepackage[letterpaper,hmargin=1.02in,vmargin=.84in]{geometry}
\usepackage[hidelinks]{hyperref}
\newtheorem{theorem}{Theorem}
\newcommand{\dist}[1]{\left\lVert #1\right\rVert}
\title{A sharp phase neighbourhood for consecutive lonely runners}
\author{\href{https://orcid.org/0009-0006-1237-7763}{Chengyi Fang}}
\date{Revised 27 November 2026}
\begin{document}
\maketitle
\begin{abstract}
For the speeds $1,\ldots,p-1$, with $p\ge7$ prime, we determine the
loneliness as a function of all initial phases on the largest centred
cube supporting the local formula. It equals $1/p$ plus the maximum
of the absolute values of $(p-1)/2$ explicit linear forms. The first
competing peaks arise at Farey denominator sum $p+1$. An obstruction
on doubling orbits yields a larger tube in which loneliness cannot
fall below $1/p$ whenever doubling modulo sign has odd order.
\end{abstract}

\section{The phase function}
Write $\dist{x}=\min_{m\in\mathbb Z}|x-m|$. For phases
$u=(u_1,\ldots,u_{p-1})\in\mathbb R^{p-1}$ put
\[
 F(u)=\max_{t\in\mathbb R/\mathbb Z}\min_{1\le j<p}\dist{jt+u_j},
 \qquad
 M(u)=\max_{1\le a<p/2}\left|\frac{(p-a)u_a-au_{p-a}}p\right|.
\]
The maximum exists by compactness. Blanco, Criado and Santos
determine the minimum over all phases for several consecutive speed
sets, including $9/67<1/7$ for $1,\ldots,6$
\cite[Prop.~5.12(2)]{bcs}. Zhang discusses the global minimum over both
velocities and phases and asks for its minimizers \cite[Question~8.5]{zhang}.
Poliakova constructs configurations with small gaps for large runner
counts \cite{pol}. Beck and Everett give necessary short integer
relations among the velocities when a shifted gap is small
\cite[Prop.~4 and Thm.~5]{be}.

Here the velocities are fixed and the question is how the gap depends on
the phases near the common start. In the discussion of Zhang's problem,
\texttt{aadrianvv} identified the Farey tight points at zero phases and
the complementary-speed peaks with opposite orientations. Perturbing
a complementary pair raises one exposed peak as it lowers the other
\cite{aadrianvv}. We determine the full local formula, its largest
centred cube, and a separate bound on a wider tube.

\begin{theorem}\label{main}
Let $p\ge7$ be prime. Take $Q=(p-1)/2$ if $p\equiv3\pmod4$ and
$Q=(p-3)/2$ otherwise. Set $r_0=1/(2Q(p+1))$. On
$\max_j|u_j|\le r_0$ one has $F(u)=1/p+M(u)$.
Every larger centred cube contains a point where $F(u)>1/p+M(u)$.
\end{theorem}

In particular the common start is a local minimum. Within this cube
its level set is the linear section $u_a/a=u_{p-a}/(p-a)$ for each
complementary pair. This allows a different common time shift for
each pair. Here $M$ uses the indicated real lifts of the phases. Both $F$ and the linear forms are unchanged by replacing
$u_j$ by $u_j+j\tau$ for a common $\tau$.

\begin{proof}
Write $d=1/p$, $r=\max_j|u_j|\le r_0$. Choose $a$ attaining $M$,
put $b=p-a$, and take $k$ with $ka\equiv1\pmod p$. At
$t=\pm k/p-(u_a+u_b)/p$, the two phase corrections relative to
$\pm k/p$ are $X=(bu_a-au_b)/p$ and $-X$. Take the sign making
the two distances $d+|X|$ (here $d+r<1/2$). All the remaining
speeds had distance at least $2d$ before the correction and change
by at most $(1+2(p-1)/p)r$. Since $4r_0\le d$, they still have
distance at least $d+r$. This proves the lower bound.

Decompose $[0,1]$ using the reduced fractions of denominator at most
$p-1$. Consecutive terms $P/q<R/h$ satisfy $qR-hP=1$ and $q+h\ge p$,
by the standard Farey neighbour identities \cite[Ch.~III]{hw}.
For $t$ in this interval put $A=qt-P+u_q$ and $B=R-ht-u_h$.
If either is negative its absolute value is at most $r$.
Otherwise the distance to the nearer of these two runners is at most
\[
 \min(A,B)\le\frac{hA+qB}{q+h}
 =\frac{1+hu_q-qu_h}{q+h}.
\]
When $q+h=p$ the last quantity is at most $d+M$; when
$q+h\ge p+2$ it is at most $1/(p+2)+r_0\le d$.

Suppose $q+h=p+1$, and now denote the smaller denominator by $q$.
Regardless of endpoint order the numerator of the phase contribution
is bounded by $h|u_q|+qr$. One has $q\le Q$. Indeed the two
denominators are coprime, and if $p=4m+1$ the candidate $2m$ is even.
Using $M\ge((p-q)|u_q|-qr)/p$, we obtain
\[
 \frac{h|u_q|+qr}{p+1}-M
 \le \frac{q|u_q|}{p(p+1)}+
 qr\left(\frac1p+\frac1{p+1}\right)\le\frac{2qr}{p}.
\]
Thus the peak lies below or at $1/(p+1)+2Qr_0/p+M=d+M$.

For sharpness fix $q=Q$ and $h=p+1-q$. We have
$\gcd(q,p+1)=1$: the odd number $q$ differs from $(p+1)/2$ by one
or two. Give indices $q,p-q$ phase $r$, index $h$ phase $-r$, and index
$p-h$ phase $-(p-h)r/h$. The four indices are distinct; set
all other phases to zero. Then $M=(p-2q)r/p$. At $t=k/(p+1)$,
$kq\equiv1\pmod{p+1}$, the two distances for $q,h$ are
$1/(p+1)+r$, while all others are at least $2/(p+1)-r$.
For $r_0<r\le1/(2(p+1))$ the former is strictly greater than
$d+M$ and the latter at least as large as the former. Such points
occur arbitrarily close to the boundary.
\end{proof}

\section{An obstruction on a wider tube}
Let $G=(\mathbb Z/p\mathbb Z)^*/\{\pm1\}$ and
$R(u)=\min_{\tau\in\mathbb R/\mathbb Z}\max_j\dist{u_j+j\tau}$.
The order of $2$ in $G$ is odd for every $p\equiv3\pmod4$ since
$|G|=(p-1)/2$ is then odd.

\begin{theorem}\label{tube}
Let $p\ge7$ be prime. If the order of $2$ in $G$ is odd, then
\[
 R(u)\le \frac{2(p+1)}{3p-1}\,\frac1p
 \quad\Longrightarrow\quad F(u)\ge1/p .
\]
\end{theorem}
\begin{proof}
Translate time, take real phase representatives, and put
$r=\max_j|u_j|$ equal to the attained minimum $R(u)$.
Suppose $F(u)<d$. Even with lifts of magnitude at most $d$, at the
two times $\pm k/p$, $ka\equiv1\pmod p$, only the speeds $a,p-a$
can have distance strictly below $d$. At positive $k/p$ this requires
$u_a<0$ or $u_{p-a}>0$; at negative $k/p$ the signs reverse.
Consequently both phases in the pair have the same strict sign.
For any $|z|\le(d-r)/(p-1)$ the phases $u_j+jz$ still have magnitude
at most $d$. Hence no such phase vanishes, which implies
\begin{equation}\label{slack}
 |u_j|>j(d-r)/(p-1).
\end{equation}

Take $1\le a<p/2$, $h=p-a$. At $t=\pm k/p-u_h/h$ the runner
$h$ is at distance $d$, and at one sign runner $a$ is at distance
at least $d$ (its correction has size at most $2r<2d$ and $3d<1/2$).
Every correction $u_j-ju_h/h$ has magnitude at most
$(1+(p-1)/h)r\le2d$, by $h\ge(p+1)/2$ and the assumed radius.
A runner attaining distance less than $d$ must therefore have had
nominal distance exactly $2d$, so its class in $G$ is $[2a]$ and
its correction has magnitude greater than $d$.
If its phase had the same sign as $u_h$, one signed difference
would be at most $|u_j|<d$ and the other, by \eqref{slack}, would satisfy
\[
 j|u_h|/h-|u_j|<
 \frac{j}{p-1}\big[(1+(p-1)/h)r-d\big]\le d .
\]
Thus the signs attached to the classes $[a]$ and $[2a]$ must be
opposite. This is impossible around an odd cycle of doubling.
\end{proof}

The order condition is equivalent to the existence of an odd positive
integer $n$ with $2^n\equiv1$ or $-1\pmod p$. The order divides any such
return exponent, and an odd order itself is such an exponent.

For $p=59$ the exact formula holds on the cube of radius $1/3480$,
and the obstruction holds on a tube of radius $15/(22\cdot59)$.
For general $p$, determining the distance from the common-start orbit
to the first strict improvement remains open here.
For composite moduli the nonunit speeds can introduce further
active constraints, and the argument does not give the same formula.

\paragraph{Verification.}
The companion Lean development gives the exact closed-cube identity,
its sharpness, and the closed-tube bound. It checks the correspondence
between the supremum and infimum over real time and the attained
extrema on one turn. It also checks the equality between the
all-index and half-index maxima for the complementary-pair forms.
The distance to the integers is represented by $\lvert x-\operatorname{round}(x)\rvert$.
The formal proof establishes the Farey interval property by finite
mediant refinement from $0/1$ and $1/1$.
The accompanying Python file
\texttt{check.py} independently evaluates the complete piecewise-linear
envelope for rational examples using integer arithmetic. At a maximum,
opposite active slopes occur, including a cusp of a single runner.
Consequently $(i+j)t+u_i+u_j$ is integral for some $i,j$ (possibly
equal), and enumeration on one period suffices. These tests are
supplementary.

\paragraph{Contributions.}
Chengyi Fang initiated and funded MathLab, set its research goals,
supervised the work, and arranged dissemination. OpenAI Codex
coordinated the research, proof engineering, and verification.
GPT-6.2 Sol subagents assisted with mathematical auditing and publication
work. GPT-6.2 Astra supplied the $p=59$ precursor and a formalization
draft for the wider tube, and carried out independent model review.
AI systems produced the mathematics, programs, and manuscript.
Chengyi Fang has not independently checked any part of the proofs.
The original proof and exact examples accompany
\href{https://github.com/F1Justin/MathLab/tree/v1.0.0/papers/phases}{MathLab v1.0.0}.
\begin{thebibliography}{9}
\bibitem{bcs} M. Blanco, F. Criado and F. Santos,
\emph{Coloopless zonotopes and counterexamples to the Shifted Lonely
Runner Conjecture}, \href{https://arxiv.org/abs/2603.24784v4}{arXiv:2603.24784v4}.
\bibitem{pol} D. Poliakova, \emph{More (shifted) runners, less loneliness},
\href{https://arxiv.org/abs/2609.23952}{arXiv:2609.23952}.
\bibitem{zhang} Y. Zhang, \emph{Tight instances of the Lonely Runner Conjecture:
one-entry modifications and the growth of speeds},
\href{https://arxiv.org/abs/2608.13599v3}{arXiv:2608.13599v3}, Appendix A.
\bibitem{be} M. Beck and S. Everett, \emph{Lonely Runner Relations},
\href{https://arxiv.org/abs/2609.06259v3}{arXiv:2609.06259v3}.
\bibitem{aadrianvv} \texttt{aadrianvv},
\href{https://github.com/1337isnot1337/solved-conjectures/issues/110\#issuecomment-7049825074}{Farey tight points at zero phases} and
\href{https://github.com/1337isnot1337/solved-conjectures/issues/110\#issuecomment-7051219544}{opposite-orientation pair obstruction},
\emph{solved-conjectures}, issue 110, 20 November 2026.
\bibitem{hw} G. H. Hardy and E. M. Wright, \emph{An Introduction to
the Theory of Numbers}, 6th ed., Oxford University Press, 2008.
\end{thebibliography}
\end{document}
